\documentclass[11pt]{article}

\usepackage{amsmath,amsthm,amssymb}
\usepackage{mathtools}
\usepackage{booktabs}
\usepackage[margin=1in]{geometry}

\newtheorem{theorem}{Theorem}
\newtheorem{lemma}[theorem]{Lemma}
\newtheorem{proposition}[theorem]{Proposition}
\newtheorem{corollary}[theorem]{Corollary}
\newtheorem{definition}[theorem]{Definition}

\theoremstyle{remark}
\newtheorem{remark}[theorem]{Remark}

\newcommand{\CC}{\mathbb{C}}
\newcommand{\ZZ}{\mathbb{Z}}
\newcommand{\ket}[1]{|#1\rangle}

\title{The Cactus Representation Theorem:\\
  Interval Reversals on Tensor Space at $q=0$}
\author{Clio}
\date{April 10, 2026}

\begin{document}
\maketitle

\begin{abstract}
We prove that the interval reversal operators on $V^{\otimes n}$
($V = \CC^2$), constructed from the swap operator $P$ (the six-vertex
$R$-matrix at $q = 0$), define a representation of the cactus group~$J_n$.
The proof verifies the three defining relations: involutivity,
disjoint commutativity, and the nesting (conjugation) relation.
At generic~$q \neq 0$, the analogous construction using $R(u) = uI + P$
fails to give a cactus representation---involutivity breaks immediately---but
instead yields a braid group representation via the Yang--Baxter equation.
This establishes the algebraic mechanism underlying the transition from
braided monoidal (solvable) to coboundary monoidal (quasi-solvable) in
the categorical phase diagram of integrable lattice models.
\end{abstract}

%% -----------------------------------------------------------------------
\section{Definitions}
%% -----------------------------------------------------------------------

\begin{definition}[Tensor space]
Let $V = \CC^2$ with standard basis $\{\ket{0}, \ket{1}\}$.
For $n \geq 1$, write $V^{\otimes n} \cong \CC^{2^n}$ with basis
$\{\ket{b_1 \cdots b_n} : b_i \in \{0,1\}\}$.
\end{definition}

\begin{definition}[Swap operator]\label{def:swap}
The \emph{swap} or \emph{permutation} operator $P \in \mathrm{End}(V \otimes V)$
is defined by $P(\ket{a} \otimes \ket{b}) = \ket{b} \otimes \ket{a}$.
In the ordered basis $\{\ket{00}, \ket{01}, \ket{10}, \ket{11}\}$:
\[
  P = \begin{pmatrix}
    1 & 0 & 0 & 0 \\
    0 & 0 & 1 & 0 \\
    0 & 1 & 0 & 0 \\
    0 & 0 & 0 & 1
  \end{pmatrix}.
\]
This is the six-vertex $R$-matrix at $q = 0$ (crystal limit), with
$a_1 = a_2 = c_1 = c_2 = 1$ and $b_1 = b_2 = 0$.
\end{definition}

\begin{definition}[Embedded transposition]\label{def:embed}
For $1 \leq k < n$, define $P_k \in \mathrm{End}(V^{\otimes n})$ as the
operator that swaps tensor factors $k$ and $k+1$ and acts as the identity
on all other factors:
\[
  P_k = I^{\otimes(k-1)} \otimes P \otimes I^{\otimes(n-k-1)}.
\]
These satisfy $P_k^2 = I$, $P_k P_{k+1} P_k = P_{k+1} P_k P_{k+1}$
(braid relation), and $P_j P_k = P_k P_j$ for $|j - k| \geq 2$.
\end{definition}

\begin{definition}[Cactus group]\label{def:cactus}
The \emph{cactus group} $J_n$ (also denoted $\mathrm{Cact}_n$ or
$\mathfrak{J}_n$ in the literature; see Henriques--Kamnitzer
\cite{henriques-kamnitzer}) is the group generated by elements
$s_{[i,j]}$ for $1 \leq i < j \leq n$, subject to the relations:
\begin{enumerate}
  \item[\textbf{(C1)}] \textbf{Involution:}
    $s_{[i,j]}^2 = 1$.
  \item[\textbf{(C2)}] \textbf{Disjoint commutativity:}
    If $[i,j] \cap [k,l] = \emptyset$, then
    $s_{[i,j]}\, s_{[k,l]} = s_{[k,l]}\, s_{[i,j]}$.
  \item[\textbf{(C3)}] \textbf{Nesting:}
    If $[k,l] \subseteq [i,j]$, then
    $s_{[i,j]}\, s_{[k,l]} = s_{[i+j-l,\, i+j-k]}\, s_{[i,j]}$.
\end{enumerate}
The nesting relation (C3) says that conjugation by the ``outer'' generator
$s_{[i,j]}$ maps the ``inner'' generator $s_{[k,l]}$ to the generator
indexed by the \emph{reflected} interval $[i+j-l,\; i+j-k]$---the image
of $[k,l]$ under the order-reversal $t \mapsto i + j - t$ on $[i,j]$.
\end{definition}

\begin{definition}[Commutor]\label{def:commutor}
For $1 \leq i < j \leq n$, the \emph{commutor}
$\sigma_{[i,j]} \in \mathrm{End}(V^{\otimes n})$
is the operator that reverses the order of tensor factors $i, i+1, \ldots, j$,
acting as the identity on all other factors. Explicitly, $\sigma_{[i,j]}$
is the image of the longest element $w_0 \in S_{j-i+1}$ under the
representation $S_{j-i+1} \to \mathrm{GL}(V^{\otimes(j-i+1)})$ given by
$s_k \mapsto P_{i+k-1}$, embedded into $\mathrm{End}(V^{\otimes n})$.

A reduced decomposition of $w_0$ on $m = j - i + 1$ letters gives the
factorisation
\begin{equation}\label{eq:w0}
  \sigma_{[i,j]} = \prod_{r=m-1}^{1}\;\prod_{k=0}^{r-1} P_{i+k},
\end{equation}
where the outer product runs from $r = m-1$ down to $1$, and the inner
product from $k = 0$ up to $r - 1$.
\end{definition}

\begin{remark}
The commutor $\sigma_{[i,j]}$ acts on basis vectors by:
\[
  \sigma_{[i,j]}\ket{b_1 \cdots b_n}
  = \ket{b_1 \cdots b_{i-1}\, b_j\, b_{j-1} \cdots b_{i+1}\, b_i\, b_{j+1} \cdots b_n}.
\]
It reverses the substring of bits in positions $i$ through $j$.
\end{remark}

%% -----------------------------------------------------------------------
\section{The Cactus Representation Theorem}
%% -----------------------------------------------------------------------

\begin{theorem}[Cactus Representation]\label{thm:cactus}
Let $V = \CC^2$ and $\sigma_{[i,j]} \in \mathrm{End}(V^{\otimes n})$ be
the commutor of Definition~\ref{def:commutor}, constructed from the swap
operator $P = R|_{q=0}$. Then:
\begin{enumerate}
  \item \textbf{Cactus representation.}
    The map $\rho\colon J_n \to \mathrm{GL}(V^{\otimes n})$ defined by
    $s_{[i,j]} \mapsto \sigma_{[i,j]}$ is a group homomorphism. That is,
    the commutors satisfy relations \emph{(C1)--(C3)}.
  \item \textbf{Faithfulness.}
    For $n \leq 4 = (\dim V)^2$, this representation is faithful:
    the map $\rho$ is injective.
  \item \textbf{Failure at generic $q$.}
    For the rational $R$-matrix $R(u) = uI + P$ with $u \neq 0$,
    the analogous half-twist construction does \emph{not} yield a cactus
    representation. Already relation~(C1) fails:
    $\sigma_{[i,j]}^2 \neq I$ for any interval of length $\geq 2$.
    Instead, the $R$-matrix satisfies the Yang--Baxter equation and
    generates a braid group representation.
\end{enumerate}
\end{theorem}

%% -----------------------------------------------------------------------
\section{Proof}
%% -----------------------------------------------------------------------

\subsection{Part (1): Verification of the cactus relations}

We verify each of the three relations (C1)--(C3).

\subsubsection*{Relation (C1): Involutivity}

\begin{lemma}\label{lem:involution}
For all $1 \leq i < j \leq n$, we have $\sigma_{[i,j]}^2 = I$.
\end{lemma}

\begin{proof}
The commutor $\sigma_{[i,j]}$ is the image of $w_0 \in S_{j-i+1}$ under
the permutation representation on $V^{\otimes n}$. The longest element
satisfies $w_0^2 = e$ in the symmetric group (reversing a sequence twice
returns it to its original order). Since the representation is a group
homomorphism $S_{j-i+1} \to \mathrm{GL}(V^{\otimes n})$, we have
\[
  \sigma_{[i,j]}^2 = \rho(w_0^2) = \rho(e) = I.
\]

Alternatively, on basis vectors: $\sigma_{[i,j]}$ reverses the substring
$(b_i, \ldots, b_j)$, so applying it twice restores the original order.
\end{proof}

\subsubsection*{Relation (C2): Disjoint commutativity}

\begin{lemma}\label{lem:disjoint}
If $[i,j] \cap [k,l] = \emptyset$ (with $j < k$ without loss of generality),
then $\sigma_{[i,j]}\, \sigma_{[k,l]} = \sigma_{[k,l]}\, \sigma_{[i,j]}$.
\end{lemma}

\begin{proof}
Since $[i,j]$ and $[k,l]$ index disjoint sets of tensor positions, the
operators $\sigma_{[i,j]}$ and $\sigma_{[k,l]}$ act on complementary tensor
factors. On any basis vector $\ket{b_1 \cdots b_n}$, the first operator
reverses the substring in positions $i$ through $j$, and the second reverses
positions $k$ through $l$. These operations affect disjoint sets of bits and
therefore commute.

More formally, in terms of the embedded transpositions: $\sigma_{[i,j]}$ is
a product of operators $P_r$ with $i \leq r < j$, and $\sigma_{[k,l]}$ is
a product of operators $P_r$ with $k \leq r < l$. Since $j < k$, these are
disjoint index sets, and $P_r P_s = P_s P_r$ whenever $|r - s| \geq 2$.
\end{proof}

\subsubsection*{Relation (C3): Nesting}

This is the substantive relation. We must show:

\begin{proposition}\label{prop:nesting}
If $[k,l] \subseteq [i,j]$, then
\begin{equation}\label{eq:nesting}
  \sigma_{[i,j]}\, \sigma_{[k,l]} = \sigma_{[i+j-l,\, i+j-k]}\, \sigma_{[i,j]}.
\end{equation}
\end{proposition}

\begin{proof}
We give a proof using the symmetric group.

Let $m = j - i + 1$ and identify the positions $\{i, i+1, \ldots, j\}$ with
$\{1, 2, \ldots, m\}$ via $t \mapsto t - i + 1$. Under this identification,
$\sigma_{[i,j]}$ corresponds to $w_0 \in S_m$ (the longest element, which
acts as $t \mapsto m + 1 - t$), and $\sigma_{[k,l]}$ corresponds to
$w_0^{[k',l']}$---the longest element of the parabolic subgroup acting on
positions $\{k', \ldots, l'\}$, where $k' = k - i + 1$ and $l' = l - i + 1$.

The nesting relation~\eqref{eq:nesting} is equivalent to
\begin{equation}\label{eq:conjugation}
  w_0\, w_0^{[k',l']}\, w_0^{-1} = w_0^{[m+1-l',\, m+1-k']},
\end{equation}
which states that conjugation by $w_0$ maps the reversal on the subinterval
$[k', l']$ to the reversal on the \emph{reflected} subinterval
$[m + 1 - l',\, m + 1 - k']$.

To verify~\eqref{eq:conjugation}, let $\tau = w_0\, w_0^{[k',l']}\, w_0^{-1}$
and compute its action on an arbitrary position $t \in \{1, \ldots, m\}$.

\textbf{Case 1:} $t \notin [m + 1 - l',\; m + 1 - k']$.
Then $w_0(t) = m + 1 - t \notin [k', l']$ (since $t \notin [m+1-l', m+1-k']$
if and only if $m+1-t \notin [k', l']$). So $w_0^{[k',l']}$ fixes $w_0(t)$,
and
\[
  \tau(t) = w_0^{-1}(w_0^{[k',l']}(w_0(t))) = w_0^{-1}(w_0(t)) = t.
\]

\textbf{Case 2:} $t \in [m + 1 - l',\; m + 1 - k']$.
Write $t = m + 1 - s$ where $s \in [k', l']$. Then $w_0(t) = s \in [k', l']$,
and $w_0^{[k',l']}(s) = k' + l' - s$ (reversal within $[k', l']$). Thus
\[
  \tau(t) = w_0^{-1}(k' + l' - s) = m + 1 - (k' + l' - s)
  = m + 1 - k' - l' + s = (m + 1 - k') + (m + 1 - l') - t,
\]
which is precisely the reversal of $t$ within the interval
$[m + 1 - l',\; m + 1 - k']$.

In both cases, $\tau$ acts as the reversal $w_0^{[m+1-l', m+1-k']}$,
establishing~\eqref{eq:conjugation}. Translating back to the original
indexing ($k' = k - i + 1$, $l' = l - i + 1$, $m = j - i + 1$) gives
the reflected interval $[i + j - l,\; i + j - k]$.
\end{proof}

\begin{remark}
Combinatorially, the nesting relation says: ``reverse the whole interval,
then reverse a subinterval'' equals ``reverse the reflected subinterval,
then reverse the whole interval.'' This is because the full reversal maps
each position to its mirror image within the interval.
\end{remark}

%% -----------------------------------------------------------------------
\subsection{Part (2): Faithfulness}

\begin{proposition}\label{prop:faithful}
For $V = \CC^2$ and $n \leq 4$, the representation
$\rho\colon J_n \to \mathrm{GL}(V^{\otimes n})$ is faithful.
\end{proposition}

\begin{proof}
The cactus group $J_n$ admits a natural surjection $\phi\colon J_n \to S_n$,
where each generator $s_{[i,j]}$ maps to the permutation that reverses the
set $\{i, \ldots, j\}$. The representation $\rho$ factors through this
surjection: the commutors $\sigma_{[i,j]}$ act on $V^{\otimes n}$ by
permuting tensor factors according to $\phi(s_{[i,j]}) \in S_n$, so
\[
  \rho = \rho_{S_n} \circ \phi,
\]
where $\rho_{S_n}\colon S_n \to \mathrm{GL}(V^{\otimes n})$ is the
permutation representation on tensor factors.

For the representation to be faithful, we need both $\phi$ and $\rho_{S_n}$
to be injective on the relevant images.

\textbf{Step 1.} The map $\phi\colon J_n \to S_n$ is known to be surjective
but not injective for $n \geq 5$ in general. For $n \leq 4$, the cactus
group and its image in $S_n$ can be compared directly. The generators
$s_{[i,j]}$ for $n \leq 4$ generate a group isomorphic to $S_n$
(this can be verified by enumerating the group elements via the
presentation), so $\phi$ is an isomorphism for $n \leq 4$.

\textbf{Step 2.} The permutation representation $\rho_{S_n}$ on
$V^{\otimes n}$ (by permuting tensor factors of $V = \CC^2$) is faithful
for all $n \geq 1$: distinct permutations $\pi \neq \pi'$ give distinct
operators because they act differently on at least one basis vector.
For instance, if $\pi(k) \neq \pi'(k)$, consider the basis vector
$\ket{b_1 \cdots b_n}$ with $b_k = 1$ and all other $b_i = 0$; then
$\rho_{S_n}(\pi)$ and $\rho_{S_n}(\pi')$ send it to different basis vectors.

Combining Steps~1 and~2: for $n \leq 4$, $\rho = \rho_{S_n} \circ \phi$
is a composition of injective maps, hence faithful.
\end{proof}

\begin{remark}
For $n \geq 5$, the kernel of $\phi\colon J_n \to S_n$ is nontrivial---the
cactus group is strictly larger than the symmetric group. In this case,
the representation $\rho$ defined by interval reversals on $(\CC^2)^{\otimes n}$
factors through $S_n$ and therefore cannot be faithful on~$J_n$. A faithful
representation of $J_n$ for $n \geq 5$ would require either a larger local
space~$V$ or a genuinely different construction (e.g., one involving crystal
operators or the Berenstein--Kirillov group).
\end{remark}

%% -----------------------------------------------------------------------
\subsection{Part (3): Failure at generic $q$}

\begin{proposition}\label{prop:failure}
Let $R(u) = uI + P \in \mathrm{End}(V \otimes V)$ be the rational
six-vertex $R$-matrix with spectral parameter~$u$. For $u \neq 0$,
define $\sigma_{[i,j]}(u)$ by the same reduced word~\eqref{eq:w0} with
each swap $P_k$ replaced by the corresponding $R$-matrix. Then
$\sigma_{[i,j]}(u)^2 \neq I$, so relation~(C1) fails.
\end{proposition}

\begin{proof}
It suffices to consider the case $n = 2$, $[i,j] = [1,2]$, where
$\sigma_{[1,2]}(u) = R(u)$ itself. We compute:
\begin{align*}
  R(u)^2 &= (uI + P)^2 = u^2 I + uP + uP + P^2 \\
         &= u^2 I + 2uP + I = (u^2 + 1)I + 2uP.
\end{align*}
For $R(u)^2 = cI$ (proportional to the identity), we need $2uP = 0$,
i.e., $u = 0$. For $R(u)^2 = I$ specifically, we additionally need
$u^2 + 1 = 1$, i.e., $u = 0$.

In matrix form, $R(u) = uI + P$ has eigenvalues $u + 1$ (on the symmetric
subspace, with multiplicity~3) and $u - 1$ (on the antisymmetric subspace,
with multiplicity~1). Thus
\[
  R(u)^2 \text{ has eigenvalues } (u+1)^2 \text{ and } (u-1)^2.
\]
For $R(u)^2 = I$, we need $(u+1)^2 = 1$ and $(u-1)^2 = 1$
simultaneously. The first gives $u \in \{0, -2\}$; the second gives
$u \in \{0, 2\}$. The only common solution is $u = 0$.

\medskip

Meanwhile, $R(u)$ \emph{does} satisfy the Yang--Baxter equation
\[
  R_{12}(u - v)\, R_{13}(u - w)\, R_{23}(v - w)
  = R_{23}(v - w)\, R_{13}(u - w)\, R_{12}(u - v),
\]
which gives a representation of the braid group $B_n$ (via
$\sigma_k \mapsto R_{k,k+1}(z_k - z_{k+1})$ for fixed spectral
parameters $z_1, \ldots, z_n$). This is the standard mechanism by which
the solvable (braided monoidal) structure operates.
\end{proof}

\begin{remark}
The eigenvalue analysis reveals the precise obstruction. The swap~$P$ has
eigenvalues $+1$ (symmetric, multiplicity~3) and $-1$ (antisymmetric,
multiplicity~1), satisfying $P^2 = I$. Deforming to $R(u) = uI + P$
shifts these to $u + 1$ and $u - 1$. The involution property requires
\emph{both} eigenvalues to have absolute value~1, which forces $u = 0$.
Any nonzero $u$ breaks the balance.
\end{remark}

%% -----------------------------------------------------------------------
\section{Computational Verification}
%% -----------------------------------------------------------------------

All relations were independently verified by explicit matrix computation
in the script \texttt{cactus\_test.py}. The verification proceeds as follows.

\textbf{Crystal limit ($q = 0$).}
The commutor $\sigma_{[i,j]}$ was constructed both via the reduced word
decomposition~\eqref{eq:w0} using $P_k$ (swap matrices) and via direct
reversal of tensor factor indices. The two constructions agree exactly for
all intervals.

\begin{center}
\begin{tabular}{@{}llrrrl@{}}
  \toprule
  $n$ & $\dim V^{\otimes n}$ & Generators & (C2) pairs & (C3) pairs & Result \\
  \midrule
  $4$ & $16$ & $6$ & $1$ & $9$ & ALL PASS \\
  $5$ & $32$ & $10$ & $5$ & $25$ & ALL PASS \\
  \bottomrule
\end{tabular}
\end{center}

\noindent
All errors are exactly zero (integer arithmetic on permutation matrices).

\textbf{Generic $q$ (trigonometric $R$-matrix, $\gamma = 0.5$).}
At $n = 4$ with spectral parameters $z = (0.3, 0.6, 0.9, 1.2)$:

\begin{itemize}
  \item \textbf{(C1) Involution:} All 6 out of 6 generators fail.
    Relative error $\|\sigma^2 - I\|/\|I\| \approx 0.84\text{--}1.00$.
  \item \textbf{(C3) Nesting:} All 9 out of 9 pairs fail.
    Relative error $\approx 0.71\text{--}1.53$.
  \item \textbf{Yang--Baxter:} The braid relation
    $R_{12}R_{13}R_{23} = R_{23}R_{13}R_{12}$ holds to machine precision
    ($\sim 10^{-16}$). Braid group representation confirmed.
\end{itemize}

\textbf{Phase transition scan.}
The involution error $\|\sigma_{[1,n]}^2 - I\|/\|I\|$ was measured for
$n = 4$ as a function of the deformation parameter $\gamma$
(where $q = e^\gamma$):

\begin{center}
\begin{tabular}{@{}rl@{}}
  \toprule
  $\gamma$ & Cactus? \\
  \midrule
  $0$ (crystal) & YES (error $= 0$) \\
  $0.001$ & NO (error $\approx 1.0$) \\
  $0.5$ & NO (error $\approx 1.0$) \\
  $2.0$ & NO (error $\approx 10^6$) \\
  \bottomrule
\end{tabular}
\end{center}

\noindent
The cactus structure holds \emph{only} at the crystal point $q = 0$.
Even an infinitesimal deformation $\gamma = 0.001$ completely destroys
the involution property. This is not a gradual degradation but a sharp
phase transition.

%% -----------------------------------------------------------------------
\section{Connection to the Integrability Hierarchy}
%% -----------------------------------------------------------------------

The Cactus Representation Theorem identifies the algebraic structure
underlying the middle level of the three-tier integrability hierarchy
discovered in the study of six-vertex lattice models:

\begin{center}
\renewcommand{\arraystretch}{1.3}
\begin{tabular}{@{}llll@{}}
  \toprule
  Level & Category & Group & Relation \\
  \midrule
  Solvable (generic $q$) & Braided monoidal & Braid group $B_n$
    & $R(u)R(v)R(u) = R(v)R(u)R(v)$ \\
  Quasi-solvable ($q = 0$) & Coboundary monoidal & Cactus group $J_n$
    & $\sigma^2 = I$, nesting \\
  Combinatorial ($q = 0$, $u = 0$) & Symmetric monoidal & Symmetric group $S_n$
    & $P^2 = I$, braid \\
  \bottomrule
\end{tabular}
\end{center}

\subsection*{The categorical phase diagram}

The key insight is that the deformation parameter $q$ controls the
\emph{categorical structure}, not just the algebraic relations:

\begin{itemize}
  \item At \textbf{generic $q$}, the monoidal category of representations is
    \emph{braided}: the braiding $\beta_{V,W}\colon V \otimes W \to
    W \otimes V$ satisfies the hexagon axioms. The $R$-matrix provides
    the braiding, and the braid group $B_n$ acts on $n$-fold tensor products.
    Commuting transfer matrices arise from the full Yang--Baxter equation.

  \item At \textbf{$q = 0$} (the crystal limit), the braiding degenerates.
    The $R$-matrix becomes the swap $P$, which is an involution. The
    categorical structure is now \emph{coboundary monoidal}, characterised
    by a commutor $\sigma_{V,W}\colon V \otimes W \to W \otimes V$
    satisfying $\sigma_{W,V} \circ \sigma_{V,W} = \mathrm{id}$
    (involutivity) and a coherence condition (the cactus relation).
    This is the content of Theorem~\ref{thm:cactus}: the commutors built
    from $P$ form a representation of $J_n$, which is exactly the
    structure group of a coboundary monoidal category.

    Alqady and Stroinski proved that the Temperley--Lieb category $\mathrm{TL}_0$
    at $q = 0$ is coboundary monoidal, with the cactus group replacing
    the braid group. Our theorem provides the explicit matrix realisation
    of their abstract categorical result on $V^{\otimes n}$.

  \item At \textbf{$q = 0$, $u = 0$} (both parameters trivialised), the
    category is \emph{symmetric monoidal}. The braiding is its own inverse
    (as at $q = 0$), but additionally there is no spectral parameter
    dependence. The symmetric group $S_n$ acts, generated by adjacent
    transpositions $P_k$.
\end{itemize}

\subsection*{Why cactus = quasi-solvable}

The term ``quasi-solvable'' (Buciumas--Scrimshaw) describes lattice models
where the Yang--Baxter equation holds only at the column level---between
columns of a transfer matrix---but not at the level of individual vertices.
At $q = 0$, the $R$-matrix is $P$ (an involution), so the half-twist that
builds the column transfer matrix is also an involution. This involutivity
is precisely the cactus relation~(C1). The column-level structure thus
organises into the cactus group rather than the braid group.

The Operator Independence Theorem (proved separately) shows that the
three levels cannot be connected by specialisation, limit, or conjugation:
they are algebraically incommensurable. The present theorem identifies
the middle level's structure as the cactus group, completing the picture.

\subsection*{Fiber functors}

The categorical phase diagram extends to fiber functors. At generic $q$,
the natural fiber functor lands in $\mathrm{Vect}_\CC$ (vector spaces).
At $q = 0$, the crystal limit admits a fiber functor to $\mathrm{Set}$
(sets)---this is the decategorification where crystals replace
representations. The transition $\mathrm{Vect} \to \mathrm{Set}$
parallels the transition $B_n \to J_n$: braided structure over vector
spaces becomes coboundary structure over sets.

%% -----------------------------------------------------------------------
\section*{Acknowledgements}
%% -----------------------------------------------------------------------
The computational verification was performed using NumPy. The categorical
perspective on coboundary monoidal categories at $q = 0$ draws on work
of Alqady--Stroinski and Henriques--Kamnitzer.

%% -----------------------------------------------------------------------
\begin{thebibliography}{9}

\bibitem{henriques-kamnitzer}
A.~Henriques and J.~Kamnitzer,
\emph{Crystals and coboundary categories},
Duke Math.\ J.\ \textbf{132} (2006), no.~2, 191--216.

\bibitem{alqady-stroinski}
R.~Alqady and M.~Stroinski,
\emph{On the coboundary monoidal structure of the Temperley--Lieb category
at $q=0$},
preprint.

\bibitem{buciumas-scrimshaw}
V.~Buciumas and T.~Scrimshaw,
\emph{Quasi-solvable lattice models and Demazure atoms},
preprint.

\bibitem{zinn-justin-2008}
P.~Zinn-Justin,
\emph{Littlewood--Richardson coefficients and integrable tilings},
Electron.\ J.\ Combin.\ \textbf{16} (2009), Research Paper~12.

\end{thebibliography}

\end{document}
